\documentclass[11pt]{article}
\usepackage[T1]{fontenc}
\usepackage[utf8]{inputenc}
\usepackage{mathpazo}
\usepackage{microtype}
\usepackage[margin=1.25in]{geometry}
\usepackage[hidelinks]{hyperref}
\setlength{\parindent}{1.4em}
\setlength{\parskip}{0pt}
\linespread{1.07}
\title{\Large\bfseries Hypernormal Cinema\\[0.5ex]\normalsize Individual Editions and the Limits of Correction\\[0.5ex]\normalsize Outline of the Monograph}
\author{Flyxion\\[0.3ex]\small Independent Researcher}
\date{November 2026}
\begin{document}
\maketitle

The monograph is organized around one question that the opening chapter makes concrete: how much variation can a production sustain while remaining open to correction. Each chapter after the first grows from a problem that the lawn example exposed, and the order follows the order in which those problems forced themselves. The chapters fall into five parts. The first two supply what the opening chapter assumed, namely a measure of perceptual difference and an account of what a scene commits to. The third and fourth parts follow the life of an edition, from its structure through correction, allocation, release, and the viewer's lifetime. The last part treats exclusion, generation, and cases.

\section*{Part I. Foundations of Variation}

\paragraph{Chapter 1. The Space of Possible Films.} Written. The exact count, the separation of enumeration, recognition, and realization, the lawn scene through three productions, counting by transfer matrices, footprints of an amendment, collision-free repair, the scope of uniqueness, the ledger, and the lifetime repertoire.

\paragraph{Chapter 2. Perceptual Equivalence and the Measurement of Difference.} Written. The opening chapter assumes that perceptual classes exist and that two editions are equivalent exactly when they share a class in every slot. This chapter replaces the assumption with a treatment of dissimilarity. It separates detection, comprehension, and affective response, each with its own pseudometric, and states what an experiment would have to measure for each. It develops the failure of transitivity for thresholded similarity into packing and covering numbers, shows how the sustainable audience depends on the choice between them, and defines the comparison across local productions that the worldwide promise requires. The chapter proves what the formalism permits and marks the points at which evidence about viewers must replace mathematics.

\paragraph{Chapter 3. Histories, Folds, and What a Scene Commits To.} Written. The verdict of Chapter 1 locates the difference between an error in the film and an act by a character in the history layer, and the history as defined records only what happened. This chapter annotates histories with dramatic function, so that a departure from competence can be marked as intended, and states the guarantee that a correction cannot delete evidence on which a later reinterpretation depends. The engine scene, in which an ordinary repair is later understood as diligence, negligence, or concealment, is the governing example. The chapter also treats noise overcapture, extending the quotient of histories by future equivalence from what can happen to what can be shown, so that a generator tracks a wet plate only when it matters to an action, a later shot, or an inference.

\section*{Part II. The Structure of an Edition}

\paragraph{Chapter 4. Slots, Takes, and the Graph of Editions.} Written. The lawn scene has constraints only between neighbors and at single slots. This chapter generalizes the layered graph to a constraint structure with relations between arbitrary slots, and tests the generalization on two cases. When the constraints form a chain, the new counting method must recover the lawn totals of Chapter 1 exactly. It must then handle a constraint between two nonadjacent shots, namely which hand holds a glass before and after an insert. The example defines the relevant notions of width, the frontier width of the scope graph and the carry width given by future equivalence of prefixes, and shows how width sets the cost of counting and the memory for unranking while the per-viewer cost of unranking stays small. The chapter also records that counting admissible editions is intractable for unrestricted constraint structures.

\paragraph{Chapter 5. Continuity Beyond Neighbors.} Written. A glass held in one hand across a cut, a strip of grass that stays mown for the rest of the scene, and a line of dialogue that runs through an insert are constraints that span distance. The earlier volume treats agreement on overlaps by sheaf conditions. This chapter uses them to compute the footprint of a correction in general, shows how a correction derived from one violated overlap finds every overlap it implicates, and gives propagation procedures with bounds on how far a repair can travel. The containment lemma of Chapter 1 is the special case in which a neutral take stops the chain.

\paragraph{Chapter 6. Sound, Dialogue, and Time.} Written. Sound occupies time in a way that image does not, and a repair that fits the picture can break the soundtrack. The chapter takes the acoustic occupation analysis of an existing film as a second worked case, with measured data, and extends the model of editions to continuous timing constraints. It treats recovery windows, the tolerance of cuts to timing shifts, and the asymmetry of those tolerances that the earlier volume's saturating response model already shows.

\section*{Part III. Correction}

\paragraph{Chapter 7. Reports and Adjudication.} Written. The three-step verdict of Chapter 1, comprising the viewer's claim, the footage the ledger shows was assigned, and the script history, is developed into a full account of reports. It separates the edition assigned from the edition actually played. An assignment ledger can establish the former, while a report about a superseded cut may require a revision identifier reported by playback to establish the latter. The chapter covers a taxonomy of defects, reports that describe the footage mistakenly, aggregation of many reports about one slot, and the standard of proof for a report that would supply evidence for a new rule. Ledger retention becomes a requirement of adjudication, since a verdict about a past edition can be reached only if the take that was played can be identified and recovered.

\paragraph{Chapter 8. Amendments, Footprints, and Correction Events.} Written. An upheld report has two possible outcomes. It can reveal a violation of a rule already in force, in which case the defect lay in the assembly or in the assigned take and the rule set does not change. Or it can supply evidence for an amendment, in which case a new rule must be derived from the instance. The chapter treats how a rule is derived from an upheld report, how its footprint is computed and its general form recovered from the specific instance, how the atomic correction event of Chapter 1 is formed, and when it cannot be formed because the repairs would break a promise. Replacement footage, relaxation of a rule that the history permits, and the order in which corrections arrive are treated as parts of one process.

\paragraph{Chapter 9. Capacity and Robustness.} Written. The book's central question has two precise answers, and this chapter gives the first. Capacity asks whether the corrected production has enough distinct editions to serve its audience from scratch, as the number of perceptual editions or the packing number, a function of rules and footage. Chapter 1 shows that capacity decreases when constraints are added to a fixed pool. The chapter develops robustness to families of anticipated amendments, which asks what reserve a production needs in order to survive them, bounds on the capacity lost by a correction in terms of its footprint, and the conditions under which a production with zero spare capacity at the premiere can be corrected at all. Capacity does not say that the editions can be reached from the current assignments. That is the subject of the next chapter.

\paragraph{Chapter 10. Allocation Under Change.} Written. The second precise answer concerns repairability. It asks whether the corrected editions can be reached from the current assignments within the permitted disturbance. A production can have a hundred distinguishable editions for sixty viewers and still fail to repair two of them, if both have the same sole acceptable replacement. Capacity is therefore necessary and not sufficient. The chapter develops the Hall condition for collision-free repair with survivors held fixed, its generalization to minimum-cost assignment in which valid editions may move at a price, online allocation as corrections arrive, the dependence of one viewer's edition on all the others, and reproducibility from a ledger of bounded size together with a retained archive. It asks whether the reserve that robustness requires is usable from a particular allocation.

\section*{Part IV. Release and Repertoire}

\paragraph{Chapter 11. Locales and the Worldwide Premiere.} Written. A release has one premiere instant and many local productions, each with its own cast, language, and setting. The criterion of uniqueness carried over from Chapter 2 decides what composes. If uniqueness means exact inequivalence, local allocation composes when the productions' class sets are disjoint. If it means perceptual separation by more than a tolerance, disjoint classes alone are insufficient, because independently selected editions from different locales must also satisfy the distance threshold across locales. The concrete case is a sequence that uses shared footage, in which the local cast, language, and setting leave no perceptible trace. The chapter treats the shared specification of the story, the passages at which editions of different productions may coincide, and the minimum footage a locale must produce to serve its audience.

\paragraph{Chapter 12. Ledgers, Archives, and the Release Slider.} Written. The ledger records assignments, and recovering an old edition requires the assets and assembly information they reference. The chapter treats retention policy, the cost of retaining superseded takes, the reconstruction of an edition from a recorded tuple, the privacy of viewer identifiers in a ledger, and the interface semantics of the slider, including what it shows for a date before a premiere and for a scheduled premiere that has no edition yet.

\paragraph{Chapter 13. The Lifetime Repertoire.} Written. A viewer has a finite viewing life. The chapter formalizes the budget, the allocation of hours across films, revisits, and release dates, and the selection and scheduling of a repertoire from the enormous space of possible films, in which the space itself never has to be enumerated. It treats how revisits change the number of distinct films a person can see and what a repertoire has to guarantee about the editions inside it.

\paragraph{Chapter 14. Sequence and Reinterpretation.} Written. The budget of the preceding chapter treats films as items, and they are not. A later film can change what an earlier one means, and the order of release is therefore a variable. The chapter treats a repertoire as a structure on films and not as a set, asks whether releasing the same film in the tenth year of a viewer's life does something that releasing it in the first would miss, and develops the formal tools for dependence of meaning on order. It is the most open chapter of the book and is written as a program.

\section*{Part V. Boundaries and Cases}

\paragraph{Chapter 15. Exclusion and Reference Corpora.} Written. A production may declare material that its editions must not reproduce. Exact subtraction by digest is arithmetic and changes the count by the number removed. Approximate subtraction depends on a representation and a threshold, removes a neighborhood around each excluded item, and interacts with packing numbers. The chapter counts the editions that survive a declared policy, treats shared cinematic convention as reusable grammar that a policy must be able to retain while excluding specific combinations of imagery, performance, and timing, and keeps the count apart from any determination of legal status.

\paragraph{Chapter 16. Generation Under Constraint.} Written. The preceding chapters describe editions assembled from takes. This chapter asks what a generator must represent to produce takes on demand under the same rules, so that the pool is not fixed in advance. It treats admissible sampling, the cost of rejecting samples that violate the rules, how a correction is communicated to a generator, and the hierarchical direction of a generator by nested commands, in the manner of the earlier volume.

\paragraph{Chapter 17. Case Studies.} Written. Three studies carry the machinery through complete examples. The engine scene follows three local productions and a large audience through delayed recognition. The lawn and a dishwashing scene follow a correction through a full ledger. A measured passage of an existing film follows the model of acoustic occupation through a correction to its own record. Each study states what the theory predicts, what it does not determine, and what an experiment would have to settle.

\paragraph{Chapter 18. Open Problems.} Written. A consolidated list of what the book leaves undecided, ordered by what each would unlock, with the empirical questions about viewers kept apart from the questions that are purely mathematical.

\section*{Appendices}

\paragraph{Appendix A. Counting and Matching.} Written. Transfer matrices, ranking and unranking, Hall's theorem and minimum-cost assignment, packing and covering numbers, each with complete proofs for the finite case used in the text.

\paragraph{Appendix B. Verification Scripts.} Written. The scripts that reproduce every count and worked example in the book, with the instances they were run on stated plainly.

\paragraph{Appendix C. Notation and Relation to the Earlier Volume.} Written. A table of notation, the correspondence between the layers of the earlier volume and the objects of this one, and the points at which the two books use different words for the same thing.

\section*{Status}

All eighteen chapters and the three appendices are written and are merged into one book, \texttt{hypernormal-cinema.tex}. Chapters compile standalone as well, so cross-references between chapters are typed as numbers and should be checked against the merged book after any renumbering. Chapter 18 and Appendix C have no verification script, and Chapter 17 contains an invented dishwashing scene.

\end{document}
